%% Lecture 09 -- Field variation and four-current
\section{Lecture 09 -- Field Variation and Four-Current}\label{sec:lec09}

\begin{center}
\begin{tabular}{ll}
  \textbf{Date:}\quad 05/05/2026 & \\
\end{tabular}
\end{center}
\hrule\vspace{1.5em}

\subsection*{Overview}
The previous lecture constructed the electromagnetic field tensor
\begin{equation}\label{eq:lec09:F-def}
  F_{\mu\nu} = \partial_\mu A_\nu - \partial_\nu A_\mu
\end{equation}
and argued that Lorentz invariance, gauge invariance, and the superposition
principle almost uniquely determine the free-field action.  This lecture begins
the variational derivation of Maxwell's equations.  The main idea is simple:
the complete system has two kinds of degrees of freedom.  Varying the particle
worldlines gives the Lorentz force law; varying the electromagnetic field
$A_\mu(x)$ gives the field equations.

The live lecture reached the following points:
\begin{itemize}
  \item the total action for point charges plus the electromagnetic field;
  \item the first variation of the free-field action;
  \item the definition of charge density, current density, and four-current;
  \item the proof that the four-current obeys the covariant continuity equation.
\end{itemize}

The derivation of Maxwell's equations is therefore almost complete: one still
has to vary the interaction term and combine it with the free-field variation.
That combination will give the inhomogeneous Maxwell equations
$\partial_\mu F^{\mu\nu} = (4\pi/c)J^\nu$.

%------------------------------------------------------------
\subsection{The Total Action}\label{subsec:lec09-total-action}

\subsubsection*{Physical Motivation}
A charged particle and the electromagnetic field are not separate mechanical
problems.  The particle moves in the field, and the field is sourced by the
particle.  The action must therefore contain three pieces: a free-particle
piece, an interaction piece, and a free-field piece.  Each term has a clear
role: inertia, coupling, and field dynamics.

For $N$ point charges, labelled by $n=1,\ldots,N$, the total action is
\begin{equation}\label{eq:lec09-total-action}
  \boxed{
  S =
  -\sum_{n=1}^N m_n c\int ds_n
  -\frac{1}{c}\sum_{n=1}^N e_n\int_{\Gamma_n} A_\mu\,dx_n^\mu
  -\frac{1}{16\pi c}\int F_{\mu\nu}F^{\mu\nu}\,d^4x .
  }
\end{equation}
\textit{Physical significance:} The first term says that a free relativistic
particle extremizes its proper time.  The second term couples the particle
worldline $\Gamma_n$ to the four-potential.  The third term gives the
electromagnetic field its own dynamics.

Here $x^\mu=(ct,\mathbf{x})$, $d^4x=dx^0\,d^3x=c\,dt\,d^3x$, and the metric
convention is $(+,-,-,-)$.  The factor of $1/c$ in the field action is the
standard Gaussian-units normalization when the four-volume element is written
as $d^4x$.

\subsubsection*{Why the field action is quadratic}

The free-field part is not guessed arbitrarily.  It is forced by three physical
requirements:
\begin{enumerate}
  \item \textbf{Lorentz invariance:} the action must be a scalar.
  \item \textbf{Gauge invariance:} the action may depend on $A_\mu$ only through
  gauge-invariant combinations, especially $F_{\mu\nu}$.
  \item \textbf{Superposition:} the field equations must be linear, so the
  action must be quadratic in the field.
\end{enumerate}

The derivative matrix $\partial_\mu A_\nu$ can be decomposed into a symmetric
and an antisymmetric part:
\begin{equation}\label{eq:lec09-sym-antisym}
  \partial_\mu A_\nu
  =
  \frac{1}{2}\left(\partial_\mu A_\nu+\partial_\nu A_\mu\right)
  +
  \frac{1}{2}\left(\partial_\mu A_\nu-\partial_\nu A_\mu\right).
\end{equation}
The antisymmetric part is precisely $F_{\mu\nu}/2$.  Under a gauge
transformation,
\begin{equation}\label{eq:lec09-gauge}
  A_\mu \longrightarrow A_\mu+\partial_\mu f ,
\end{equation}
the antisymmetric combination is unchanged:
\begin{align}
  F_{\mu\nu}
  &\longrightarrow
  \partial_\mu(A_\nu+\partial_\nu f)
  -\partial_\nu(A_\mu+\partial_\mu f) \notag\\
  &= F_{\mu\nu}
  +\partial_\mu\partial_\nu f
  -\partial_\nu\partial_\mu f
  = F_{\mu\nu}.
  \label{eq:lec09-F-gauge-inv}
\end{align}
Thus the allowed quadratic Lorentz scalar is
\begin{equation}\label{eq:lec09-free-action}
  \boxed{
  S_{\rm field}
  =
  -\frac{1}{16\pi c}\int F_{\mu\nu}F^{\mu\nu}\,d^4x .
  }
\end{equation}
\textit{Physical significance:} Gauge symmetry throws away the symmetric part
of $\partial_\mu A_\nu$ and keeps only the curl-like object $F_{\mu\nu}$.  This
is the covariant version of the familiar fact that physical electromagnetic
fields are built from derivatives of potentials.

\subsubsection*{Why $d^4x$ is Lorentz invariant}

The four-volume element behaves like an ordinary volume element under a change
of variables:
\begin{equation}\label{eq:lec09-jacobian}
  d^4x' = \left|\det \Lambda\right|\,d^4x .
\end{equation}
For a proper Lorentz transformation, $\det\Lambda=1$, so
\begin{equation}\label{eq:lec09-d4x-scalar}
  \boxed{d^4x'=d^4x.}
\end{equation}
\textit{Physical significance:} A Lorentz boost squeezes spatial volume and
stretches time in just the compensating way needed to leave the four-dimensional
spacetime volume unchanged.

\begin{center}
\begin{tikzpicture}[scale=0.95, >=Latex]
  \draw[->] (-0.2,0) -- (5.2,0) node[right] {$x$};
  \draw[->] (0,-0.2) -- (0,4.2) node[above] {$ct$};
  \draw[thick, blue] (0.8,0.6) -- (3.6,3.4);
  \draw[thick, blue!60] (1.0,0.6) -- (3.8,3.4);
  \draw[thick, blue!60] (0.8,0.9) -- (3.6,3.7);
  \draw[thick, blue!60] (1.0,0.9) -- (3.8,3.7);
  \node[blue] at (3.9,3.15) {boosted cell};
  \draw[dashed] (0.8,0.6) -- (1.0,0.6) -- (1.0,0.9) -- (0.8,0.9) -- cycle;
  \node at (1.25,0.35) {$d^4x$};
  \node[align=center] at (2.7,1.1) {spatial contraction\\and time dilation\\compensate};
\end{tikzpicture}
\end{center}

%------------------------------------------------------------
\subsection{Varying the Free Electromagnetic Field}\label{subsec:lec09-field-variation}

\subsubsection*{Physical Motivation}
In mechanics, one varies the path $q(t)\to q(t)+\delta q(t)$ and demands that
the first-order change in the action vanish.  For a field, the same idea is
applied at every spacetime point:
\begin{equation}\label{eq:lec09-A-variation}
  A_\mu(x)\longrightarrow A_\mu(x)+\delta A_\mu(x).
\end{equation}
The variation $\delta A_\mu$ is assumed to vanish on the boundary of the
spacetime integration region, just as $\delta q(t_1)=\delta q(t_2)=0$ in
ordinary mechanics.

\subsubsection*{First variation of $F_{\mu\nu}$}

Since $F_{\mu\nu}$ is linear in $A_\mu$,
\begin{align}
  F_{\mu\nu}[A+\delta A]
  &=
  \partial_\mu(A_\nu+\delta A_\nu)
  -\partial_\nu(A_\mu+\delta A_\mu) \notag\\
  &=
  F_{\mu\nu}[A]
  +\partial_\mu\delta A_\nu
  -\partial_\nu\delta A_\mu.
  \label{eq:lec09-F-var-expanded}
\end{align}
Therefore
\begin{equation}\label{eq:lec09-delta-F}
  \boxed{
  \delta F_{\mu\nu}
  =
  \partial_\mu\delta A_\nu-\partial_\nu\delta A_\mu .
  }
\end{equation}
\textit{Physical significance:} A small change in the potential changes the
field tensor only through derivatives of that change.  This is why integration
by parts will move derivatives from $\delta A_\mu$ onto $F^{\mu\nu}$ and produce
field equations.

\subsubsection*{First variation of the free-field action}

Starting from
\begin{equation}\label{eq:lec09-Sf-start}
  S_{\rm field}
  =
  -\frac{1}{16\pi c}\int F_{\mu\nu}F^{\mu\nu}\,d^4x ,
\end{equation}
the first-order variation is
\begin{align}
  \delta S_{\rm field}
  &=
  -\frac{1}{16\pi c}\int
  \left[
    \delta F_{\mu\nu}F^{\mu\nu}
    +F_{\mu\nu}\delta F^{\mu\nu}
  \right]d^4x .
  \label{eq:lec09-dSf-two-terms}
\end{align}
The two terms are equal.  Indeed,
\begin{align}
  F_{\mu\nu}\delta F^{\mu\nu}
  &=
  F_{\mu\nu}\eta^{\mu\alpha}\eta^{\nu\beta}\delta F_{\alpha\beta}
  =
  F^{\alpha\beta}\delta F_{\alpha\beta}.
  \label{eq:lec09-raising-deltaF}
\end{align}
Renaming the dummy indices $\alpha,\beta$ back to $\mu,\nu$ gives
\begin{equation}\label{eq:lec09-two-terms-equal}
  F_{\mu\nu}\delta F^{\mu\nu}=F^{\mu\nu}\delta F_{\mu\nu}.
\end{equation}
Thus
\begin{equation}\label{eq:lec09-dSf-one-term}
  \delta S_{\rm field}
  =
  -\frac{1}{8\pi c}\int F^{\mu\nu}\delta F_{\mu\nu}\,d^4x .
\end{equation}
Substituting \eqref{eq:lec09-delta-F},
\begin{align}
  \delta S_{\rm field}
  &=
  -\frac{1}{8\pi c}\int
  F^{\mu\nu}
  \left(\partial_\mu\delta A_\nu-\partial_\nu\delta A_\mu\right)d^4x .
  \label{eq:lec09-dSf-sub}
\end{align}
The two terms again combine.  In the second term, exchange the dummy indices
$\mu\leftrightarrow\nu$:
\begin{align}
  -F^{\mu\nu}\partial_\nu\delta A_\mu
  &\longrightarrow
  -F^{\nu\mu}\partial_\mu\delta A_\nu
  =
  +F^{\mu\nu}\partial_\mu\delta A_\nu,
  \label{eq:lec09-antisym-combine}
\end{align}
where $F^{\nu\mu}=-F^{\mu\nu}$.  Therefore
\begin{equation}\label{eq:lec09-dSf-before-parts}
  \delta S_{\rm field}
  =
  -\frac{1}{4\pi c}\int
  F^{\mu\nu}\partial_\mu\delta A_\nu\,d^4x .
\end{equation}

\subsubsection*{Integration by parts in spacetime}

Use the product rule
\begin{equation}\label{eq:lec09-product-rule}
  F^{\mu\nu}\partial_\mu\delta A_\nu
  =
  \partial_\mu\!\left(F^{\mu\nu}\delta A_\nu\right)
  -(\partial_\mu F^{\mu\nu})\delta A_\nu .
\end{equation}
Then
\begin{align}
  \delta S_{\rm field}
  &=
  -\frac{1}{4\pi c}\int
  \partial_\mu\!\left(F^{\mu\nu}\delta A_\nu\right)d^4x
  +
  \frac{1}{4\pi c}\int
  (\partial_\mu F^{\mu\nu})\delta A_\nu\,d^4x .
  \label{eq:lec09-dSf-parts}
\end{align}
The first term is a four-dimensional boundary term.  Since
$\delta A_\nu=0$ on the boundary, it vanishes.  Hence
\begin{equation}\label{eq:lec09-dSf-final}
  \boxed{
  \delta S_{\rm field}
  =
  \frac{1}{4\pi c}\int
  (\partial_\mu F^{\mu\nu})\delta A_\nu\,d^4x .
  }
\end{equation}
\textit{Physical significance:} The quantity $\partial_\mu F^{\mu\nu}$ is the
field-theory analogue of ``force equals derivative of the action''.  Once the
source term is included, it will be set equal to the charge-current source.

%------------------------------------------------------------
\subsection{Charge Density and Current Density for Point Charges}\label{subsec:lec09-rho-J}

\subsubsection*{Physical Motivation}
To vary the interaction term with respect to the field, we want to express the
point particles as a spacetime source distribution.  A point charge is not a
smooth lump of charge; all of its charge is concentrated at one position.
Dirac delta functions are the correct language for that concentration.

For point charges $e_n$ moving along spatial trajectories $\mathbf{x}_n(t)$, the
charge density is
\begin{equation}\label{eq:lec09-rho-def}
  \boxed{
  \rho(t,\mathbf{x})
  =
  \sum_{n=1}^N e_n\,
  \delta^{(3)}\!\left(\mathbf{x}-\mathbf{x}_n(t)\right).
  }
\end{equation}
\textit{Physical significance:} Integrating $\rho$ over a region returns the
total charge of the particles inside that region.  If a particle is outside the
region, its delta function contributes zero.

The ordinary current density is charge density times velocity:
\begin{equation}\label{eq:lec09-J-def}
  \boxed{
  \mathbf{J}(t,\mathbf{x})
  =
  \sum_{n=1}^N e_n\,\dot{\mathbf{x}}_n(t)\,
  \delta^{(3)}\!\left(\mathbf{x}-\mathbf{x}_n(t)\right).
  }
\end{equation}
\textit{Physical significance:} $\mathbf{J}$ measures how much charge crosses a
unit area per unit time.  A static charge contributes to $\rho$ but not to
$\mathbf{J}$.

\begin{center}
\begin{tikzpicture}[scale=1.0, >=Latex]
  \draw[->] (-0.2,0) -- (5.2,0) node[right] {$x$};
  \draw[->] (0,-0.2) -- (0,3.8) node[above] {$ct$};
  \draw[thick, blue] (0.8,0.5) .. controls (1.4,1.3) and (2.0,1.5) .. (2.5,2.2)
    .. controls (3.0,2.9) and (3.7,3.0) .. (4.4,3.4);
  \fill[red] (2.5,2.2) circle (2pt);
  \draw[dashed] (2.5,2.2) -- (2.5,0);
  \draw[dashed] (2.5,2.2) -- (0,2.2);
  \node[blue] at (4.25,3.1) {$\Gamma_n$};
  \node[red] at (3.25,2.25) {$\mathbf{x}_n(t)$};
  \node at (2.5,-0.3) {$\mathbf{x}_n$};
  \node[left] at (0,2.2) {$ct$};
  \node[align=center] at (2.7,0.9) {delta function\\selects the\\worldline};
\end{tikzpicture}
\end{center}

%------------------------------------------------------------
\subsection{The Four-Current}\label{subsec:lec09-four-current}

\subsubsection*{Physical Motivation}
The electromagnetic field is written covariantly using four-vectors and tensors.
The source must also be written covariantly.  The natural object is the
four-current: its time component is charge density, and its spatial components
are ordinary current density.

Define
\begin{equation}\label{eq:lec09-four-current-components}
  \boxed{
  J^\mu=(c\rho,\mathbf{J}).
  }
\end{equation}
\textit{Physical significance:} Charge density and current density are not
separate relativistic objects.  They are frame-dependent components of one
four-vector, just as energy and momentum are components of one four-momentum.

Using \eqref{eq:lec09-rho-def} and \eqref{eq:lec09-J-def}, this can be written
componentwise as
\begin{equation}\label{eq:lec09-four-current-dt}
  \boxed{
  J^\mu(x)
  =
  \sum_{n=1}^N e_n\,
  \frac{dx_n^\mu}{dt}\,
  \delta^{(3)}\!\left(\mathbf{x}-\mathbf{x}_n(t)\right).
  }
\end{equation}
Indeed, for $\mu=0$,
\begin{equation}\label{eq:lec09-J0-check}
  J^0
  =
  \sum_{n=1}^N e_n\,\frac{d(ct)}{dt}\,
  \delta^{(3)}\!\left(\mathbf{x}-\mathbf{x}_n(t)\right)
  =
  c\rho,
\end{equation}
and for $\mu=i$,
\begin{equation}\label{eq:lec09-Ji-check}
  J^i
  =
  \sum_{n=1}^N e_n\,\frac{dx_n^i}{dt}\,
  \delta^{(3)}\!\left(\mathbf{x}-\mathbf{x}_n(t)\right)
  =
  J^i_{\rm spatial}.
\end{equation}

\subsubsection*{Why this is really a four-vector}

At first glance \eqref{eq:lec09-four-current-dt} is not manifestly covariant
because it contains $dt$ and a three-dimensional delta function.  The more
covariant way to see the structure is to insert a delta function in time as
well, so that the source is supported on the entire worldline:
\begin{equation}\label{eq:lec09-current-covariant}
  \boxed{
  J^\mu(x)
  =
  c\sum_{n=1}^N e_n\int d\tau_n\,
  u_n^\mu(\tau_n)\,
  \delta^{(4)}\!\left(x-x_n(\tau_n)\right),
  }
\end{equation}
where
\begin{equation}\label{eq:lec09-four-velocity}
  u_n^\mu=\frac{dx_n^\mu}{d\tau_n}
\end{equation}
is the four-velocity of particle $n$.  The four-dimensional delta function is
defined by
\begin{equation}\label{eq:lec09-delta4-def}
  \int d^4x\,\delta^{(4)}(x-a)=1.
\end{equation}
Since $d^4x$ is Lorentz invariant, $\delta^{(4)}(x-a)$ transforms as a scalar
under proper Lorentz transformations.  The charge $e_n$ is a scalar, $d\tau_n$
is a scalar, and $u_n^\mu$ is a four-vector.  Therefore $J^\mu$ is a
four-vector.

\begin{equation}\label{eq:lec09-J-vector-result}
  \boxed{J^\mu \text{ is a Lorentz four-vector.}}
\end{equation}
\textit{Physical significance:} Different observers disagree on how much source
is called ``charge density'' and how much is called ``current'', but they agree
that both are components of the same spacetime current.

%------------------------------------------------------------
\subsection{Charge Conservation as a Continuity Equation}\label{subsec:lec09-continuity}

\subsubsection*{Physical Motivation}
Charge cannot disappear from one place without flowing somewhere else.  Locally,
this means that a decrease of charge density inside a small volume must be
balanced by current flowing through its boundary.  The continuity equation is
the differential statement of this local conservation law.

Start with the ordinary current density:
\begin{equation}\label{eq:lec09-Ji-start}
  J^i(t,\mathbf{x})
  =
  \sum_{n=1}^N e_n\,\dot{x}_n^i(t)\,
  \delta^{(3)}\!\left(\mathbf{x}-\mathbf{x}_n(t)\right).
\end{equation}
Its spatial divergence is
\begin{align}
  \partial_i J^i
  &=
  \sum_{n=1}^N e_n\,\dot{x}_n^i(t)\,
  \partial_i
  \delta^{(3)}\!\left(\mathbf{x}-\mathbf{x}_n(t)\right).
  \label{eq:lec09-divJ-first}
\end{align}
Use the delta-function identity
\begin{equation}\label{eq:lec09-delta-identity}
  \partial_i
  \delta^{(3)}\!\left(\mathbf{x}-\mathbf{x}_n(t)\right)
  =
  -\frac{\partial}{\partial x_n^i}
  \delta^{(3)}\!\left(\mathbf{x}-\mathbf{x}_n(t)\right).
\end{equation}
Then
\begin{align}
  \partial_i J^i
  &=
  -\sum_{n=1}^N e_n\,\dot{x}_n^i(t)\,
  \frac{\partial}{\partial x_n^i}
  \delta^{(3)}\!\left(\mathbf{x}-\mathbf{x}_n(t)\right) \notag\\
  &=
  -\frac{\partial}{\partial t}
  \sum_{n=1}^N e_n\,
  \delta^{(3)}\!\left(\mathbf{x}-\mathbf{x}_n(t)\right) \notag\\
  &=
  -\frac{\partial \rho}{\partial t}.
  \label{eq:lec09-divJ-final}
\end{align}
Thus
\begin{equation}\label{eq:lec09-continuity-3d}
  \boxed{
  \frac{\partial \rho}{\partial t}+\nabla\cdot\mathbf{J}=0.
  }
\end{equation}
\textit{Physical significance:} Charge conservation is local.  If $\rho$ changes
inside a region, the change is accounted for by current crossing the boundary,
not by charge being created or destroyed.

\subsubsection*{Covariant form}

Since $x^0=ct$ and $J^0=c\rho$,
\begin{equation}\label{eq:lec09-time-part}
  \partial_0 J^0
  =
  \frac{\partial}{\partial x^0}(c\rho)
  =
  \frac{1}{c}\frac{\partial}{\partial t}(c\rho)
  =
  \frac{\partial\rho}{\partial t}.
\end{equation}
Therefore the ordinary continuity equation
\eqref{eq:lec09-continuity-3d} becomes
\begin{equation}\label{eq:lec09-continuity-covariant}
  \boxed{
  \partial_\mu J^\mu=0.
  }
\end{equation}
\textit{Physical significance:} Charge conservation is Lorentz covariant.  It
is not a statement true only in one preferred frame; it is the vanishing
four-divergence of the source four-vector.

%------------------------------------------------------------
\subsection{How This Leads to Maxwell's Equations}\label{subsec:lec09-next-step}

\subsubsection*{Physical Motivation}
The free-field variation knows how the electromagnetic field wants to evolve in
empty space.  The interaction term tells the field where its sources are.
Maxwell's equations are obtained by demanding that the total variation vanish
for every allowed $\delta A_\mu$.

Using the four-current, the interaction term can be rewritten as a spacetime
integral:
\begin{equation}\label{eq:lec09-Sint-current}
  \boxed{
  S_{\rm int}
  =
  -\frac{1}{c^2}\int J^\mu A_\mu\,d^4x .
  }
\end{equation}
\textit{Physical significance:} The compact contraction $J^\mu A_\mu$ is the
covariant version of ``sources couple to potentials''.  The charge density
couples to the scalar potential, and the current density couples to the vector
potential.

Varying this term with respect to $A_\mu$ gives
\begin{equation}\label{eq:lec09-dSint}
  \delta S_{\rm int}
  =
  -\frac{1}{c^2}\int J^\nu\delta A_\nu\,d^4x .
\end{equation}
Combining with \eqref{eq:lec09-dSf-final} and imposing
$\delta S_{\rm field}+\delta S_{\rm int}=0$ for arbitrary $\delta A_\nu$ gives
\begin{equation}\label{eq:lec09-maxwell-covariant}
  \boxed{
  \partial_\mu F^{\mu\nu}
  =
  \frac{4\pi}{c}J^\nu .
  }
\end{equation}
\textit{Physical significance:} This is the covariant form of the inhomogeneous
Maxwell equations.  The $\nu=0$ component is Gauss's law, and the spatial
components are the Ampere--Maxwell law.

Although the live lecture stopped just before completing this last combination,
all ingredients were prepared: the field variation \eqref{eq:lec09-dSf-final},
the source $J^\mu$, and the conservation law $\partial_\mu J^\mu=0$.

\subsection*{Summary}

\begin{itemize}
  \item The complete action contains free-particle, interaction, and free-field
  terms.
  \item Gauge invariance selects the antisymmetric field tensor
  $F_{\mu\nu}=\partial_\mu A_\nu-\partial_\nu A_\mu$.
  \item The first variation of the free-field action is
  $\delta S_{\rm field}=(4\pi c)^{-1}\int(\partial_\mu F^{\mu\nu})\delta A_\nu\,d^4x$.
  \item Point charges are represented by delta-function charge and current
  densities.
  \item The four-current $J^\mu=(c\rho,\mathbf{J})$ is a Lorentz four-vector.
  \item Local charge conservation is the covariant equation
  $\partial_\mu J^\mu=0$.
  \item Combining the field variation with the interaction variation gives
  $\partial_\mu F^{\mu\nu}=(4\pi/c)J^\nu$, the inhomogeneous Maxwell equations.
\end{itemize}
